% ============================================================
%  ch03d：§3.6 定理 A 的证明、§3.7 玩具模型（原文 pp.22–25）
%  边界：本块起于 §3.6 标题（p022 下半页），止于推论 3.7.2.1（p025 中部）；
%  其后的 §4 归 ch04a。p022 上半页的定理 3.5.2.1 归 ch03c。
% ============================================================
\subsection{定理 A 的证明}

我们现在由定理 3.5.2.1 与先前的观察推出本文的主定理。此处提醒，我们已经假设 $G$ 是半单的。\footnote{这个假设使讨论略有简化。适当表述后，本文的结果对一般约化群同样适用。}% [Tutor Note: 原文句末在脚注标记之后残留 "here."（"我们已经假设 G 是半单的。\footnote{...} here."），语序不通，疑为作者编辑残留；译文把 here 的含义并入「此处」。]

\pnum{首先，我们按照 [22] 与 [30] 回忆 Langlands 参数的显式含义。}

设 $\sigma$ 是 $\LS^{\arthm}_{\bG}$ 的一个 $\e$-点，即 $X$ 上的一个 Weil $\bG$-局部系统。我们得到一个映射
\[
  \mathrm{ev}_\sigma : \Exc_{\bG} \to \e
\]
它把函数 $f\in\Exc=\Gamma(\LS^{\arthm}_{\bG},\OO)$ 送到其在 $\sigma$ 处的值。我们滥用记号，也用 $\mathrm{ev}_\sigma$ 记诱导映射（通过取 $H^0$ 得到）$H^0(\Exc)\to\e$，它是经典交换代数之间的映射。我们用 $\mathfrak m_\sigma\subseteq H^0(\Exc)$ 记相应的极大理想。

现在回忆（来自 [4]，它建立在 [22] 与 [30] 之上）：$\Exc$——从而 $H^0(\Exc)$——作用在 $\Aut^{\unr}_{G,c}$ 上。

于是我们定义
\[
  \Aut^{\unr}_{G,c,[\sigma]} \subseteq \Aut^{\unr}_{G,c}
\]
为右端的 $\mathfrak m_\sigma$-挠部分，即 $\psi\in\Aut^{\unr}_{G,c,[\sigma]}$ 当且仅当 $\mathfrak m_\sigma^n\cdot\psi=0$ 对 $n\gg 0$ 成立。

我们令 $\Aut^{\unr}_{G,\cusp,[\sigma]} := \Aut^{\unr}_{G,c,[\sigma]}\cap\Aut^{\unr}_{G,\cusp}$。

\begin{warning}
由于 $\Aut^{\unr}_{G,\cusp}$ 是有限维的\footnote{我们提醒（这里用到 $G$ 半单），存在定义在 $\F_q$ 上的拟紧开子集 $\mathcal U\subseteq\Bun_G$，使得任何 $\psi\in\Aut^{\unr}_{G,\cusp}$ 在 $\mathcal U(\F_q)$ 之外为零（参见 [9] 命题 1.4.6 的层论版本）；由于 $\mathcal U(\F_q)$ 有限，断言显然成立。}，$\Aut^{\unr}_{G,\cusp}$ 分解为直和：
\begin{equation}
  \Aut^{\unr}_{G,\cusp} \simeq \bigoplus_{\sigma/\sim}\Aut^{\unr}_{G,\cusp,[\sigma]}.
\end{equation}
（这里隐含的等价关系 $\sim$ 把 $\mathfrak m_{\sigma_1}=\mathfrak m_{\sigma_2}$ 的 $\sigma_1$ 与 $\sigma_2$ 等同起来；根据 [22] 命题 0.38，这恰好发生在 $\sigma_1$ 与 $\sigma_2$ 具有等价半单化之时。）
\end{warning}

然而，对紧支撑自守函数，我们并没有类似的分解 (3.13)；存在这样的函数，它们不属于任何 $\Aut^{\unr}_{G,c,[\sigma]}$。

\pnum{回忆\footnote{参见 [26] II.2.4 中一个强得多的论断。} $\mathrm{ps}\text{-}\Eis(\Aut^{\unr}_{T,c})\cap\Aut^{\unr}_{G,\cusp}=0$。}

因此，只需证明任何 $\psi\in\Aut^{\unr}_{G,c,[\triv]}\subseteq\Aut^{\unr}_{G,c}$ 都落在映射 $\mathrm{ps}\text{-}\Eis:\Aut^{\unr}_{T,c}\to\Aut^{\unr}_{G,c}$ 的像中。这将是我们的目标。

\pnum{回忆 $\delta_{\bG}\in\Exc_{\bG}:=\Gamma(\LS^{\arthm}_{\bG},\OO)$。由于 $\Aut^{\unr}_{G,c}$ 被 $\Exc_{\bG}$ 作用，我们可以反转 $\delta_{\bG}$ 的作用：}
\[
  \Aut^{\unr}_{G,c}[\delta_{\bG}^{-1}] := \colim_{\delta_{\bG}\cdot-}\Aut^{\unr}_{G,c}.
\]

现在我们借助受限几何朗兰兹（以及命题 3.4.3.1）从谱 Eisenstein 级数转译过来。由引理 3.5.1.1 (1)，映射 $\mathrm{ps}\text{-}\Eis$ 把 $\delta_{\bT}$ 与 $\delta_{\bG}$ 分别在 $\Aut^{\unr}_{T,c}$ 与 $\Aut^{\unr}_{G,c}$ 上的作用衔接起来。此外，诱导映射
\[
  \mathrm{ps}\text{-}\Eis : \Aut^{\unr}_{T,c}[\delta_{\bT}^{-1}]
  \longrightarrow \Aut^{\unr}_{G,c}[\delta_{\bG}^{-1}]
\]
由定理 3.5.2.1 是满射。

这意味着，对我们的给定\footnote{需要说明：这对任何紧支撑自守函数都成立，但若该形式另有 Langlands 参数，则该论断可能基本上是空的。}自守函数 $\psi$，存在整数 $n\ge 0$，使得 $\delta_{\bG}^n\cdot\psi=\mathrm{ps}\text{-}\Eis(\psi_0)$，其中 $\psi_0\in\Aut^{\unr}_{T,c}$。

令 $\lambda\in\e$ 为 $\delta_{\bG}$ 在平凡局部系统 $\triv\in\LS^{\arthm}_{\bG}$ 处的值。注意 $(\delta_{\bG}-\lambda)\in\mathfrak m_{\triv}\subseteq H^0(\Exc)$，故对 $m\gg 0$ 有
\[
  (\delta_{\bG}-\lambda)^m\cdot\psi=0.
\]

依假设，$\lambda\ne 0$。因此可以找到多项式 $q(t)\in\e[t]$ 使
\[
  q(t)\cdot t^n \equiv 1 \pmod{(t-\lambda)^m}.
\]

于是我们显然得到
\[
  \mathrm{ps}\text{-}\Eis(q(\delta_{\bT})\cdot\psi_0)
  = q(\delta_{\bG})\cdot\mathrm{ps}\text{-}\Eis(\psi_0)
  = q(\delta_{\bG})\cdot\delta_{\bG}^n\cdot\psi
  = \psi.
\]

这就完成了论证。

\subsection{定理 3.5.2.1 的玩具模型}

我们现在给出一个更简单的设定，在其中定理 3.5.2.1 的某种形式成立。我们最终会把定理 3.5.2.1 的证明归约到这个特殊情形。我们所考虑的特殊情形，是关于 Grothendieck--Springer 消解的一个标准结果。

\pnumt{类比。}{}作为类比，我们把图表
\[
\begin{tikzcd}[column sep=huge, row sep=large]
  & \LS^{\restr}_{\bB} \arrow[dl, "q"{above}] \arrow[dr, "p"{above}] & \\
  \LS^{\restr}_{\bT} & & \LS^{\restr}_{\bG}
\end{tikzcd}
\]
换成图表：
\begin{equation}
\begin{tikzcd}[column sep=huge, row sep=large]
  & \mathbf{B}\check B \arrow[dl] \arrow[dr] & \\
  \mathbf{B}\check T & & \mathbf{B}\check G.
\end{tikzcd}
\end{equation}

代替 Frobenius，我们考虑 (3.14) 中的每一项连同其恒等自同态。取该映射之下的不动点，我们得到图表：
\[
\begin{tikzcd}[column sep=huge, row sep=large]
  & \check B^{\mathrm{ad}}/\check B \arrow[dl] \arrow[dr] & \\
  \check T^{\mathrm{ad}}/\check T & & \check G^{\mathrm{ad}}/\check G
\end{tikzcd}
\]

这里对代数群 $H$，$H^{\mathrm{ad}}/H$ 表示 $H$ 通过共轭作用在自身上所得的（叠）商；我们提醒\footnote{这里的同构存在一些符号含糊。通常这无关紧要，就当前的讨论而言亦然。但它在本文后面会有影响。我们在 §4.4.2 中澄清隐含的符号约定。}，$H^{\mathrm{ad}}/H\simeq(\mathbf{B}H)^{S^1}=(\mathbf{B}H)^{\id=\id}$。我们注意到，自然映射 $\OO_{H^{\mathrm{ad}}/H}\to\omega_{H^{\mathrm{ad}}/H}$ 是同构。

因此，§3.2 的形式体系给出典范映射：
\begin{equation}
  \Eissoyt : \Gamma\bigl(\check T^{\mathrm{ad}}/\check T,\OO_{\check T^{\mathrm{ad}}/\check T}\bigr)
  \longrightarrow
  \Gamma\bigl(\check G^{\mathrm{ad}}/\check G,\OO_{\check G^{\mathrm{ad}}/\check G}\bigr).
\end{equation}

\begin{lem}
(3.15) 的左端与右端都集中在上同调次数 $0$，且映射 $\Eissoyt$ 是满射。
\end{lem}

\begin{proof}
最直接的证明如下。我们把 $\Rep(\check T)$ 与 $\bigoplus_{\lambda\in\bar\Lambda}\Vect$ 等同，把 $\Rep(\check G)$ 与 $\bigoplus_{\lambda\in\bar\Lambda^+}\Vect$ 等同；这里 $\oplus$ 表示 $\DGCat_{\mathrm{cont}}$ 上的余积，并且我们已经默认选取了各不可约表示的同构类的代表。在此等同下，我们得到典范同构：
\begin{equation}
\begin{aligned}
  \bigoplus_{\lambda\in\bar\Lambda}k
  &\simeq \tr_{\Rep(\check T)}(\id)\,
    \bigl(\simeq\Gamma(\check T^{\mathrm{ad}}/\check T,\OO_{\check T^{\mathrm{ad}}/\check T})\bigr),\\
  \bigoplus_{\lambda\in\bar\Lambda^+}k
  &\simeq \tr_{\Rep(\check G)}(\id)\,
    \bigl(\simeq\Gamma(\check G^{\mathrm{ad}}/\check G,\OO_{\check G^{\mathrm{ad}}/\check G})\bigr).
\end{aligned}
\end{equation}

这些恒等式显然蕴涵 (3.15) 的两端都集中在次数 $0$。我们用 $e_{\check T,\bar\lambda}$（$\bar\lambda\in\bar\Lambda$）与 $e_{\check G,\bar\lambda}$（$\bar\lambda\in\bar\Lambda^+$）记由所显示的这些同构给出的向量空间的基向量。由 Borel--Weil--Bott，对支配的 $\bar\lambda$，映射 $\Rep(\check T)\to\Rep(\check G)$ 把 $\ell^{w_0(\bar\lambda)}\in\Rep(\check T)^{\vee}$（对应于 $w_0(\bar\lambda)$ 的一维表示）送到 $V^{\lambda}\in\Rep(\check G)$（最高权为 $\lambda$ 的表示）% [Tutor Note: 原文此处写作 V^λ、最高权写作 λ，与本段其余地方的 λ̄∈Λ̄ 记号不一致；译文照录。]
，从而把 $e_{\check T,w_0(\bar\lambda)}$ 送到 $e_{\check G,\bar\lambda}$；由此得到满射性。
\end{proof}

\begin{remark}
我们注意到（尽管我们并不需要这一点），在 (3.16) 中，同构 $\bigoplus_{\lambda\in\bar\Lambda^+}k\simeq\Gamma(\check G^{\mathrm{ad}}/\check G,\OO_{\check G^{\mathrm{ad}}/\check G})$ 把 $e_{\check G,\bar\lambda}$（记号同前）送到与 $\check G$ 的表示 $V^{\lambda}$ 相应的迹函数，对 $\check T$ 类似。因此，复合
\[
  \bigoplus_{\lambda\in\bar\Lambda}k
  \simeq \Gamma(\check T^{\mathrm{ad}}/\check T,\OO_{\check T^{\mathrm{ad}}/\check T})
  \xrightarrow{\ \Eissoyt\ }
  \Gamma(\check G^{\mathrm{ad}}/\check G,\OO_{\check G^{\mathrm{ad}}/\check G})
  \simeq \mathrm{Fun}(\check T)^{W}
  = \Bigl(\bigoplus_{\lambda\in\bar\Lambda}k\Bigr)^{W}
\]
可由 Weyl 特征标公式（以及 Borel--Weil--Bott）显式算出。
\end{remark}

\pnum{我们最后记录一个变体。}

注意我们有交换图：
\[
\begin{tikzcd}[column sep=huge, row sep=large]
  & \check B^{\mathrm{ad}}/\check B \arrow[dl] \arrow[dr] & \\
  \check T^{\mathrm{ad}}/\check T \arrow[dr] & & \check G^{\mathrm{ad}}/\check G \arrow[dl] \\
  & \check T/\!/W &
\end{tikzcd}
\]
其中 $\check T/\!/W=\Spec(\e[\bar\Lambda]^W)$ 是 GIT 商，左下方的映射是同义反复的（tautological），右下方的映射是标准的特征多项式映射（由该图表唯一刻画）。

因此，$\Eissoyt$ 是 $\mathrm{Fun}(\check T/\!/W)$-模的映射。引理 3.7.1.1 于是说明 $\Eissoyt$ 是 $\mathrm{Fun}(\check T/\!/W)$-模的满态射，故我们得到：

\begin{cor}
对任何 $g\in\Gamma(\check T/\!/W,\OO_{\check T/\!/W})$，映射
\[
  \Eissoyt : \Gamma\bigl(\check T^{\mathrm{ad}}/\check T,\OO_{\check T^{\mathrm{ad}}/\check T}\bigr)[g^{-1}]
  \longrightarrow
  \Gamma\bigl(\check G^{\mathrm{ad}}/\check G,\OO_{\check G^{\mathrm{ad}}/\check G}\bigr)[g^{-1}]
\]
是满射。
\end{cor}
